
(en se plaçant
dans \(H(K[T_1,\ldots,T_{n-1}])\) grâce aux identifications 1.4).

Compte tenu de 1.6.ii), \(\Phi(1+\ldots+X_nt^n)\equiv\Phi(1+\ldots+X_{n-1}t^{n-1})\ \mathrm{mod.}\ F^nH(K[X_i]_{i\geq1})\), et les \(\Phi(1+\ldots+X_nt^n)\) définissent un élément \(\psi\) de \(H(K[X_i]_{i\geq1})\) puisque c'est un groupe séparé complet~; on posera donc \(\Phi(\varphi)=\psi\).

Soient enfin \(A\in\operatorname{Ob}(\mathcal C)\), et \(\chi\in\Ghat(A)\)~; on pose \(\Phi(\chi)=H(f_\chi)(\psi)\), \(f_\chi\) étant l'homomorphisme de \(K[X_i]_{i\geq 1}\) dans \(A\) défini en 1.2. Comme on a déjà défini \(\Phi\) sur les éléments \(1+\ldots+X_nt^n=\varphi_n\), il faut vérifier que les deux définitions donnent la même chose, c'est à dire que~:
\[
  H(f_{\varphi_n})(\psi)=\Phi(1+\ldots+X_nt^n).
\]
Or pour tout \(p\), on a~: \(\psi\equiv \Phi(1+\ldots+X_pt^p)\ \mathrm{mod.}\ F^{p+1}H(K[X_i]_{i\geq 1})\). D'où pour tout \(p\geq n\)~; \(H(f_{\varphi_n})(\psi)\equiv H(f_{\varphi_n})\bigl(\Phi(1+\ldots+X_pt^p)\bigr)\ \mathrm{mod.}\ F^{p+1}H(K[X_i]_{i\geq 1})\). Utilisant les identifications de 1.4 et le recollement des \(\Phi(\varphi_m)\) pour \(m\) variable, on obtient~: \(H(f_{\varphi_n})(\Phi(\varphi_p))=\Phi(\varphi_n)\).
On a donc finalement~:
\[
  \forall p\geq n,\quad
  H(f_{\varphi_n})(\psi)\equiv \Phi(1+\ldots+X_nt^n)
  \ \mathrm{mod.}\ F^pH(K[X_i]_{i\geq 1}).
\]
Comme \(H\) est séparé, on en déduit l'égalité cherchée.

On a donc défini ainsi un système d'application \(\Phi_A:\Ghat(A)\to H(A)\)~; pour vérifier qu'elles constituent un morphisme de foncteurs (à valeurs dans \(\underline{\mathrm{Ens}}\)), il faut montrer que pour tout morphisme \(h:A\to A'\), le diagramme~:
\[
\begin{tikzcd}
\Ghat(A) \arrow[r,"\Ghat(h)"] \arrow[d,"\Phi_A"'] &
\Ghat(A') \arrow[d,"\Phi_{A'}"]\\
H(A) \arrow[r,"H(h)"] & H(A')
\end{tikzcd}
\]
